\documentclass{article}
\usepackage{amsmath,amssymb}
\newcommand{\R}{\mathbb{R}}
\newcommand{\SU}{\operatorname{SU}}
\newcommand{\SO}{\operatorname{SO}}
\newcommand{\su}{\mathfrak{su}}
\newcommand{\so}{\mathfrak{so}}
\begin{document}
\title{Frankel's rotation groups: numerical and symbolic encodings}
\author{mathphysics implementation notes}
\date{17 October 2026}
\maketitle
These are known mathematical identities implemented independently in Rust
and C++, not new results. Our source conventions follow Frankel
\cite[pp.~492--493, Eqs.~(19.1)--(19.6)]{Frankel};
the original PDF pages555--556 were directly checked.
Code uses zero-based indices: its $\epsilon(0,1,2)=1$ is the book's
$\epsilon_{123}=1$. These are Euclidean three-dimensional symbols, not the
two-component spinor or Lorentzian four-dimensional epsilon conventions.

\section{Matrices, group elements, and Lie algebra elements}
A general matrix is an unconstrained array. By contrast,
\begin{align*}
\SO(3)&=\{R\in M_3(\R):R^TR=I,\ \det R=1\},\\
\so(3)&=\{A\in M_3(\R):A^T=-A\},\\
\SU(2)&=\{U\in M_2(\mathbb C):U^\dagger U=I,\ \det U=1\}.
\end{align*}
Rotation matrices are not closed under matrix addition or arbitrary entry
mutation. Therefore our group types contain matrices rather than inherit
general matrix mutators. Floating membership is checked at a finite stated
tolerance; it is not an exact proof of these equations.

Frankel's basis is
\[
 E_1=\begin{pmatrix}0&0&0\\0&0&-1\\0&1&0\end{pmatrix},\quad
 E_2=\begin{pmatrix}0&0&1\\0&0&0\\-1&0&0\end{pmatrix},\quad
 E_3=\begin{pmatrix}0&-1&0\\1&0&0\\0&0&0\end{pmatrix}.
\]
For example,
\[
 E_1E_2=\begin{pmatrix}0&0&0\\1&0&0\\0&0&0\end{pmatrix},\quad
 E_2E_1=\begin{pmatrix}0&1&0\\0&0&0\\0&0&0\end{pmatrix},
\]
so $[E_1,E_2]=E_3$. Cyclic permutation and antisymmetry give
\[
 [E_2,E_3]=E_1,\qquad [E_3,E_1]=E_2,\qquad
 [E_i,E_j]=\epsilon_{ijk}E_k .
\]
Every implemented basis pair, including zero and reversed brackets, is tested.

\section{Epsilon, hat map, and cross product}
The alternating symbol is the sign of a permutation and vanishes on
repetitions. Its contraction identity is
\[
 \epsilon_{ijk}\epsilon_{k\ell m}
 =\delta_{i\ell}\delta_{jm}-\delta_{im}\delta_{j\ell},
 \qquad \epsilon_{ijk}\epsilon_{\ell jk}=2\delta_{i\ell}.
\]
In a general oriented coordinate frame the metric volume tensor instead has
ordered component $\sqrt{|\det g|}$; the bare symbol is not that entire tensor.

Write $\widehat v=v^iE_i$. The chosen signs give
\[
 \widehat v=\begin{pmatrix}0&-v_3&v_2\\v_3&0&-v_1\\-v_2&v_1&0\end{pmatrix},
 \qquad \widehat v\,w=v\times w .
\]
Applying the epsilon contraction to the product gives
\[
 \widehat u\,\widehat v=vu^T-(u\cdot v)I,\qquad
 [\widehat u,\widehat v]=\widehat{u\times v}.
\]
The epsilon identity also gives the Jacobi identity. Cadabra independently
reduces these structure-constant contractions; it does not infer them from
the letter $E$ without declarations.

\section{Exponential and Rodrigues formula}
Let $r^2=v\cdot v$ and $K=\widehat v$. The preceding product identity gives
\[
 K^2=vv^T-r^2I,\qquad K^3=-r^2K .
\]
Reduce each term of the matrix exponential series using this cubic relation.
The odd powers sum to $\sin r\,K/r$; the nonconstant even powers sum to
$(1-\cos r)K^2/r^2$. Thus
\[
 e^{\widehat v}=I+\frac{\sin r}{r}\widehat v+
                 \frac{1-\cos r}{r^2}\widehat v^2.
\]
The coefficients have limits1 and$1/2$ at zero. For $|n|=1$,
\[
 e^{\theta\widehat n}
 =I+\sin\theta\,\widehat n+(1-\cos\theta)\widehat n^2.
\]
Our numerical small-angle branch uses truncated coefficient series rather
than subtract nearly equal floating numbers. The symbolic branch uses the
exact ring $\mathbb Q[s,c]/(s^2+c^2-1)$, with normal form $a(s)+c\,b(s)$.
It describes one common formal angle, not a general trigonometric CAS.

\section{Pauli matrices and the normalization distinction}
Frankel gives
\[
 \sigma_1=\begin{pmatrix}0&1\\1&0\end{pmatrix},\quad
 \sigma_2=\begin{pmatrix}0&-i\\i&0\end{pmatrix},\quad
 \sigma_3=\begin{pmatrix}1&0\\0&-1\end{pmatrix}.
\]
Direct multiplication gives
\[
 \sigma_i\sigma_j=\delta_{ij}I+i\epsilon_{ijk}\sigma_k,\qquad
 [\sigma_i,\sigma_j]=2i\epsilon_{ijk}\sigma_k .
\]
The Pauli matrices are Hermitian and determinant$-1$, not elements of the
real anti-Hermitian Lie algebra
\[
 \su(2)=\{X:X^\dagger=-X,\ \operatorname{tr}X=0\}.
\]
Use explicitly distinguished bases
\[
 T_i=\frac{\sigma_i}{2i},\quad F_i=\frac{\sigma_i}{i}=2T_i,\qquad
 [T_i,T_j]=\epsilon_{ijk}T_k,\quad
 [F_i,F_j]=2\epsilon_{ijk}F_k.
\]
Consequently, coordinates in the $T$ basis have the same cross-product
bracket as $\so(3)$. Coordinates in the full $F$ basis require the factor two.
The new API names these normalizations rather than relying on letter names
used differently in older notes.

Since $(n\cdot\sigma)^2=I$ for unit $n$, separate even/odd series once more:
\[
 \exp(\theta n^iT_i)
 =\cos(\theta/2)I-i\sin(\theta/2)n\cdot\sigma ,
\]
whereas
\[
 \exp(\theta n^iF_i)
 =\cos\theta I-i\sin\theta\,n\cdot\sigma .
\]
For $n=e_3$ the latter is exactly
$\operatorname{diag}(e^{-i\theta},e^{i\theta})$.
The symbolic implementation retains whether its $s,c$ mean the full or half
angle; identical formal matrices without this metadata would be ambiguous.

\section{Quaternions, double cover, and exact symbolic charts}
For a real unit quaternion $q=(a,b)$ with $a^2+|b|^2=1$, define
\[
 U(q)=aI-i\,b\cdot\sigma,\qquad
 R(q)=(a^2-|b|^2)I+2bb^T+2a\,\widehat b .
\]
Pauli multiplication gives
\[
 (a,b)(c,d)=(ac-b\cdot d,\ ad+cb+b\times d).
\]
Independent polynomial tests verify
\[
 U(qp)=U(q)U(p),\qquad R(qp)=R(q)R(p),\qquad
 U(q)(x\cdot\sigma)U(q)^\dagger=(R(q)x)\cdot\sigma .
\]
For $q=(\cos(\theta/2),n\sin(\theta/2))$, $R(q)$ is the Rodrigues rotation.
The differential sends $T_i\mapsto E_i$ and $F_i\mapsto2E_i$.
The two lifts $q,-q$ give different SU(2) values but the same SO(3) value.
A$2\pi$ spinor turn gives$-I$, not$I$; a$4\pi$ turn gives$I$.

A useful exact symbolic parameterization is
\[
 q(t)=\frac{(1-|t|^2,\ 2t)}{1+|t|^2},\qquad t\in\R^3 .
\]
Its denominator is positive, and its unit norm follows from
$(1-|t|^2)^2+4|t|^2=(1+|t|^2)^2$.
This chart misses$-I$ in SU(2); its negative misses$+I$, so together they cover
the group. One projected chart already covers SO(3), because the two omitted/
included lifts have the same image. These are quaternion stereographic
parameters $t=n\tan(\theta/4)$, not the usual matrix Cayley parameters.
Polynomial coefficients are rational, while the formal parameters are real;
rational evaluation witnesses alone do not cover every real group value.

\begin{thebibliography}{1}
\bibitem{Frankel}
Theodore Frankel, \emph{The Geometry of Physics: An Introduction},
3rd edition, Cambridge University Press, Secs.~19.1a--b,
pp.~492--493 in the supplied PDF.
\end{thebibliography}
\end{document}
